% Figure: lift and drag on an electrodynamic (eddy-current) maglev magnet
% moving over a conducting sheet, against speed. Drag peaks at a
% characteristic speed and falls; lift rises monotonically toward an
% asymptote. Above the crossover, lift exceeds drag and the ratio grows.
\begin{figure}[htbp]
\centering
\begin{tikzpicture}[>=Latex]
\begin{axis}[
  width=12cm, height=7cm,
  xlabel={speed $v/v_c$ (characteristic speed $v_c$)},
  ylabel={force (normalised to asymptotic lift)},
  xmin=0, xmax=6,
  ymin=0, ymax=1.1,
  legend pos=south east,
  legend cell align=left,
  font=\scriptsize,
  grid=both,
  grid style={enggray!20},
]
% Let u = v/vc. Lift L(u) = u^2/(1+u^2). Drag D(u) = u/(1+u^2).
% Lift/drag = u. At u=1 lift=drag=0.5; drag peaks at u=1.
\addplot[engblue, line width=1.4pt, domain=0:6, samples=160]
  {x^2/(1+x^2)};
\addlegendentry{lift $\propto v^2/(v^2+v_c^2)$}
\addplot[engred, line width=1.2pt, dashed, domain=0:6, samples=160]
  {x/(1+x^2)};
\addlegendentry{drag $\propto v/(v^2+v_c^2)$}
% mark drag peak at u=1, value 0.5
\addplot[engamber, only marks, mark=*, mark size=2pt] coordinates {(1,0.5)};
\node[engamber, anchor=south west, font=\scriptsize] at (axis cs:1.05,0.5)
  {drag peak at $v=v_c$};
% crossover: lift=drag at u=1 too (both 0.5)
\node[enggray, anchor=west, font=\scriptsize] at (axis cs:2.6,0.2)
  {lift/drag $= v/v_c$};
\end{axis}
\end{tikzpicture}
\caption[Eddy-current maglev lift and drag]{Lift and drag on an
electrodynamic maglev magnet moving over a conducting sheet, against speed,
each normalised to the asymptotic lift. At low speed the induced eddy currents
are small and both forces vanish; drag rises to a peak at a characteristic
speed set by the sheet's resistance and the geometry, then falls as the
currents are pushed by the wave's own field to lag by a quarter cycle. Lift
rises monotonically toward an asymptote and overtakes drag above the
characteristic speed, so the lift-to-drag ratio grows in direct proportion to
speed. The vehicle is heaviest to move near the drag peak, at low speed, and
most efficient once it runs well above it, which is why an electrodynamic
maglev needs wheels to reach lift-off speed and only levitates once it is
moving fast enough to leave the drag peak behind.}
\label{fig:vol04ch08:eddy-maglev}
\end{figure}
